\section{Automated test inventory}
\label{app:tests}

% Test names and paths contain underscores; \detokenize renders them literally without manual
% escaping, and avoids the \path name clash with TikZ used elsewhere in this document.
\providecommand{\tn}[1]{\texttt{\detokenize{#1}}}

This appendix lists every automated test in the project, with the assertion each one makes and its
result. The suites run in continuous integration on every push and pull request
(\cref{sec:evaluation}). As of the run on 2026-06-29 all \textbf{71} tests pass: \textbf{62} Python
tests (executed with \texttt{pytest}) and \textbf{9} host-compiled C tests for the on-device cascade.
Line coverage of the unit-testable core (\tn{meowtion_dash/data.py} and the Cloud Function's
\tn{main.py}) is \textbf{100\,\%}; request handlers, model training, the dashboard UI and
hardware-bound firmware paths are excluded from that figure and are covered by on-hardware
verification (\cref{sec:evaluation}).

\subsection{Firmware unit tests (C, host-compiled)}

The pure collar logic is compiled and run on the host (no Zephyr, no hardware) from
\tn{firmware/test}. These cover the rules-based activity gate and the confidence-gated IMU/audio
cascade.

\begin{table}[ht]
  \centering
  \footnotesize
  \caption{Firmware C unit tests.}
  \label{tab:tests-c}
  \begin{tabularx}{\textwidth}{@{}Xr@{}}
    \toprule
    \textbf{Test and assertion} & \textbf{Result} \\
    \midrule
    \multicolumn{2}{@{}l}{\emph{\tn{test_activity.c} , rules-based activity gate}}\\
    \tn{test_motion_keeps_awake}: motion above the still threshold keeps the collar awake & Pass \\
    \tn{test_sustained_stillness_trips_rest}: sustained stillness past the hold-off enters low-power rest & Pass \\
    \tn{test_motion_resets_the_timer}: fresh motion resets the stillness timer, deferring rest & Pass \\
    \addlinespace
    \multicolumn{2}{@{}l}{\emph{\tn{test_classifier.c} , confidence-gated cascade}}\\
    \tn{test_confident_imu_skips_audio}: a confident IMU result skips the audio stage & Pass \\
    \tn{test_unsure_imu_consults_audio}: an uncertain IMU result triggers audio confirmation & Pass \\
    \tn{test_audio_kept_only_when_more_confident}: audio replaces the IMU result only when more confident & Pass \\
    \tn{test_disabled_confirm_never_runs_audio}: with confirmation disabled, audio never runs & Pass \\
    \tn{test_no_pcm_skips_audio}: with no audio samples buffered, the audio stage is skipped & Pass \\
    \tn{test_no_audio_model_skips_audio}: with no audio model loaded, the audio stage is skipped & Pass \\
    \bottomrule
  \end{tabularx}
\end{table}

\subsection{Cloud Function unit tests (Python)}

From \tn{app/firebase/functions/tests}: the model-training signal logic (\tn{test_dsp.py}) and the
request-path validation and WAV parsing (\tn{test_helpers.py}). Parametrised tests note the number of
cases they cover.

\begin{table}[ht]
  \centering
  \footnotesize
  \caption{Cloud Function Python unit tests.}
  \label{tab:tests-cloud}
  \begin{tabularx}{\textwidth}{@{}Xr@{}}
    \toprule
    \textbf{Test and assertion} & \textbf{Result} \\
    \midrule
    \multicolumn{2}{@{}l}{\emph{\tn{test_dsp.py} , windowing and normalisation}}\\
    \tn{test_per_window_norm_is_zero_mean_and_bounded}: per-window normalisation is zero-mean and bounded & Pass \\
    \tn{test_per_window_norm_matches_formula}: normalisation matches the reference formula & Pass \\
    \tn{test_per_window_norm_normalises_each_channel_independently}: each channel is normalised independently & Pass \\
    \tn{test_windows_pads_short_input_to_a_single_window}: short input is padded to a single window & Pass \\
    \tn{test_windows_slides_with_hop}: the window slides by the configured hop & Pass \\
    \tn{test_build_dataset_windows_imu_clips}: the dataset builder windows IMU clips & Pass \\
    \tn{test_build_dataset_windows_audio_clips}: the dataset builder windows audio clips & Pass \\
    \tn{test_build_dataset_skips_empty_imu}: empty IMU clips are skipped & Pass \\
    \addlinespace
    \multicolumn{2}{@{}l}{\emph{\tn{test_helpers.py} , path validation and WAV parsing}}\\
    \tn{test_safe_id_accepts_valid_segments}: \tn{_safe_id} accepts valid id segments (4 cases) & Pass \\
    \tn{test_safe_id_rejects_traversal_and_separators}: rejects path traversal, separators and spaces (10 cases) & Pass \\
    \tn{test_safe_ts_accepts_digits_only}: \tn{_safe_ts} accepts all-digit timestamps & Pass \\
    \tn{test_safe_ts_rejects_non_digits}: rejects signed, hex, spaced or traversal timestamps (7 cases) & Pass \\
    \tn{test_parse_wav_roundtrips_samples}: the WAV parser round-trips PCM samples & Pass \\
    \tn{test_parse_wav_reads_rate_from_fmt_chunk}: reads the sample rate from the fmt chunk & Pass \\
    \tn{test_parse_wav_rejects_non_wav}: rejects non-WAV input & Pass \\
    \tn{test_parse_wav_requires_a_data_chunk}: requires a data chunk & Pass \\
    \bottomrule
  \end{tabularx}
\end{table}

\subsection{Dashboard unit tests (Python)}

From \tn{app/dashboard/tests}: the data layer that turns the raw Firebase JSON into the chart-ready
frame, the time-window drill-down, the per-day segmenting used by the timeline, and the health-watch
comparison.

\begin{table}[ht]
  \centering
  \footnotesize
  \caption{Dashboard data-layer Python unit tests.}
  \label{tab:tests-dash}
  \begin{tabularx}{\textwidth}{@{}Xr@{}}
    \toprule
    \textbf{Test and assertion} & \textbf{Result} \\
    \midrule
    \multicolumn{2}{@{}l}{\emph{\tn{test_data.py}}}\\
    \tn{test_fmt_dur_under_and_over_a_minute}: duration formatting under and over a minute & Pass \\
    \tn{test_activity_dataframe_maps_class_index_to_label}: real events map the class index to the model label & Pass \\
    \tn{test_activity_dataframe_rest_event_uses_type_not_class}: the 0xFE rest sentinel falls back to the event type & Pass \\
    \tn{test_activity_dataframe_v1_event_uses_type}: legacy v1 events use the stored type & Pass \\
    \tn{test_activity_dataframe_skips_malformed_events}: events with no start time are skipped & Pass \\
    \tn{test_filter_by_activity_single_and_list}: the activity filter accepts a name or a list & Pass \\
    \tn{test_over_time_sums_duration_by_activity}: over-time aggregation sums duration by activity & Pass \\
    \addlinespace
    \multicolumn{2}{@{}l}{\emph{\tn{test_more.py}}}\\
    \tn{test_iter_cats_handles_station_nested_and_standalone}: parses station-nested and standalone cat shapes & Pass \\
    \tn{test_list_cats_dedups_and_pairs_name_id}: de-duplicates and pairs (name, id) & Pass \\
    \tn{test_period_options_day_newest_first_deduped}: day options are newest-first and de-duplicated & Pass \\
    \tn{test_filter_to_window_day_and_month}: day and month window filtering & Pass \\
    \tn{test_health_signals_recent_average_without_baseline}: recent average with no baseline yet & Pass \\
    \tn{test_fmt_time_format}: epoch-to-clock time formatting & Pass \\
    \tn{test_filter_by_weekday}: weekday filtering (single or list) & Pass \\
    \tn{test_filter_by_date_range}: date-range filtering (bounded or open-ended) & Pass \\
    \tn{test_last_n_days_keeps_recent_calendar_days}: keeps the most recent N calendar days & Pass \\
    \tn{test_health_signals_reports_drop_vs_baseline}: reports a habit drop versus baseline as a percentage & Pass \\
    \tn{test_model_labels}: reads model labels, with empty/None guards & Pass \\
    \tn{test_iter_cats_skips_non_dict_entries}: skips non-dict device and cat entries & Pass \\
    \tn{test_period_options_empty_week_and_month}: empty input and the Week/Month option paths & Pass \\
    \tn{test_filter_to_window_empty_and_week}: empty input and the 7-day window & Pass \\
    \tn{test_over_time_empty_and_day}: empty input and the hourly (Day) path & Pass \\
    \tn{test_last_n_days_and_health_signals_handle_empty}: empty-input guards & Pass \\
    \tn{test_daily_segments_rows_per_day_with_time_of_day}: events split per day with time-of-day, crossing midnight & Pass \\
    \tn{test_daily_segments_skips_zero_and_empty}: zero-duration and empty input are skipped & Pass \\
    \tn{test_trim_sparse_edge_days_drops_thin_first_and_last}: near-empty leading/trailing days are dropped & Pass \\
    \tn{test_trim_sparse_edge_days_keeps_thin_interior_day}: a thin interior day is kept & Pass \\
    \tn{test_trim_sparse_edge_days_handles_single_and_empty}: single-day and empty inputs are handled & Pass \\
    \bottomrule
  \end{tabularx}
\end{table}
